\documentclass{article}

\usepackage{graphicx}
\usepackage{amsmath}
\usepackage{booktabs}
\usepackage{float}
\usepackage{hyperref}
\usepackage[dvipsnames]{xcolor}
\usepackage{listings}
\usepackage{adjustbox}
\usepackage[
	a4paper,
	left=2.5cm,
	right=2.5cm,
	top=2.5cm,
	bottom=2.5cm,
	headheight=14pt,
	headsep=0.6cm,
	footskip=1cm
]{geometry}
\usepackage{fancyhdr}

\definecolor{nbbg}{HTML}{F7F8FA}
\definecolor{nbframe}{HTML}{D9DEE8}
\definecolor{nbkeyword}{HTML}{1F4E79}
\definecolor{nbstring}{HTML}{A31515}
\definecolor{nbcomment}{HTML}{2E8B57}
\definecolor{nbnumbers}{HTML}{6A737D}

\lstdefinestyle{nbpython}{
	language=Python,
	backgroundcolor=\color{nbbg},
	basicstyle=\ttfamily\small,
	keywordstyle=\color{nbkeyword}\bfseries,
	stringstyle=\color{nbstring},
	commentstyle=\color{nbcomment}\itshape,
	numbers=left,
	numberstyle=\tiny\color{nbnumbers},
	numbersep=8pt,
	frame=single,
	rulecolor=\color{nbframe},
	breaklines=true,
	showstringspaces=false,
	tabsize=4,
	columns=fullflexible
}

\providecommand{\tightlist}{%
	\setlength{\itemsep}{0pt}\setlength{\parskip}{0pt}%
}
\providecommand{\pandocbounded}[1]{%
	\adjustbox{max width=0.8\linewidth,max totalheight=0.65\textheight}{#1}%
}

\pagestyle{fancy}
\fancyhf{}
\fancyhead[L]{ 12-PHY-BIPGP2 General Physics Laboratory 2}
\fancyfoot[R]{ Affan Kaknjo, Andrea Pokrajčić \thepage}
\renewcommand{\headrulewidth}{0.4pt}
\renewcommand{\footrulewidth}{0.4pt}

\setlength{\parindent}{0pt}
\setlength{\parskip}{0.8em}

\fancypagestyle{plain}{
	\fancyhf{}
	\fancyhead[L]{ 12-PHY-BIPGP2 General Physics Laboratory 2}
	\fancyfoot[R]{ Affan Kaknjo, Andrea Pokrajčić \thepage}
	\renewcommand{\headrulewidth}{0.4pt}
	\renewcommand{\footrulewidth}{0.4pt}
}

\begin{document}



\section{nb2tex Repro Cases}

Synthetic notebook containing exact markdown patterns that previously
broke TeX conversion.

We arrive at values of \(R_R = 109.5\Omega\), \(R_L = 5.9\Omega\) and
\(C = 2.165 \mu F\). \ldots{} the Lissajous curve collapses into a line
45 \(\ensuremath{^\circ}\) tilted w.r.t to the x-axis.

If we let frequency \(f \to \infty\), we see that the phase shift
\(\phi(f)\) tends to \(0 \ensuremath{^\circ}\) for the RC circuit and
\(90 \ensuremath{^\circ}\) for the RL and RLC circuit.

On the other hand, if we let frequency \(f \to -\infty\), we see that
the phase shift \(\phi(f)\) tends to \(90 \ensuremath{^\circ}\) for the
RC circuit and \(0 \ensuremath{^\circ}\) for the RL and RLC circuit.

Hence it is no surprise that the series resistance
\(R = (108 \pm 5) \Omega\) acquired by the fitting procedure is almost
equal to the measured value of the resistor
resistance \(R_R=(109.5 \pm 1.095) \Omega\).

\subsection{Long Equation Split Demo}

The following equation is intentionally very long and should be
auto-split at \texttt{=} into aligned rows while preserving one equation
number in TeX output.


\begin{equation}
\begin{aligned}
Z_{RL} &= \frac{i\omega L+R_{sp}}{1+i \omega C(i\omega L+R_{sp})} \\
&= \frac{i\omega L+R_{sp}}{1 - \omega^2LC + i \omega R_{sp}C} \\
&= \frac{(i\omega L+R_{sp})(1 - \omega^2LC - i \omega R_{sp}C)}{(1 - \omega^2LC)^2 + \omega^2 (R_{sp}C)^2} \\
&= \frac{R_{sp}+\omega^2 (R_{sp}LC-R_{sp}LC) + i\left( \omega L - \omega^3L^2C - \omega R_{sp}^2C\right)}{1 + \omega^2 \left((R_{sp}C)^2- 2LC \right)+ \omega^4L^2C^2}
\end{aligned}
\label{nb2tex:eq:1}
\end{equation}


\end{document}
